Рассматривается задача обучения нейронной сети как минимизация функции потерь
\begin{equation}
\mathcal{L} : \mathbb{R}^d \to \mathbb{R},
\end{equation}
где $\theta \in \mathbb{R}^d$ обозначает вектор параметров модели. 
\emph{Ландшафтом функции потерь} называется отображение значений $\mathcal{L}(\theta)$ на параметрическое пространство, рассматриваемое как высокоразмерная невыпуклая функция.

Классическое представление ландшафта функции потерь в виде множества изолированных локальных минимумов, разделённых высокими барьерами, в общем случае не подтверждается для переобученных нейронных сетей. Вместо этого наблюдаются протяжённые области параметрического пространства, в которых функция потерь принимает близкие и малые значения.

Точка $\theta^\ast$ называется \emph{локальным минимумом}, если существует окрестность $U(\theta^\ast)$ такая, что
\begin{equation}
\mathcal{L}(\theta^\ast) \le \mathcal{L}(\theta), \quad \forall \theta \in U(\theta^\ast).
\end{equation}
При этом в высокоразмерных моделях многие точки, достигаемые в процессе обучения, могут являться седловыми точками с малыми отрицательными собственными значениями матрицы Гессе.

Для анализа связности решений используется понятие параметрического пути. 
Пусть $\theta_0, \theta_1 \in \mathbb{R}^d$. 
\emph{Путём} между ними называется непрерывное отображение
\begin{equation}
\gamma : [0,1] \to \mathbb{R}^d,
\quad
\gamma(0) = \theta_0, \ \gamma(1) = \theta_1.
\end{equation}

Наличие пути, вдоль которого значение функции потерь остаётся низким,
\begin{equation}
\mathcal{L}(\gamma(t)) \approx \max\{\mathcal{L}(\theta_0), \mathcal{L}(\theta_1)\},
\quad t \in [0,1],
\end{equation}
указывает на связность соответствующей области параметрического пространства. В литературе это явление известно как \emph{mode connectivity}, однако в настоящей работе данный термин используется лишь как описательный и не вводится как формально определённое понятие.

Для количественного описания переходов между состояниями модели вводится \emph{барьер функции потерь}. 
Пусть $\Gamma$ — фиксированный класс путей между $\theta_0$ и $\theta_1$. 
Барьер определяется как
\begin{equation}
\Delta \mathcal{L}(\theta_0, \theta_1)
=
\min_{\gamma \in \Gamma}
\max_{t \in [0,1]}
\mathcal{L}(\gamma(t))
-
\max\{\mathcal{L}(\theta_0), \mathcal{L}(\theta_1)\}.
\end{equation}
Таким образом, барьер является путь-зависимой величиной и определяется относительно выбранного класса путей.

Для описания глобальной структуры ландшафта используется понятие субуровневого множества функции потерь. 
Для фиксированного уровня $c \in \mathbb{R}$ определяется множество
\begin{equation}
\mathcal{M}_c = \{\theta \in \mathbb{R}^d \mid \mathcal{L}(\theta) \le c\}.
\end{equation}
Топологические свойства субуровневых множеств, включая их связность, определяются критическими точками функции потерь и их индексами. Связность может возникать при наличии критических точек индекса один либо в случае, когда наблюдаемые минимумы фактически являются седловыми точками.

Оптимизация функции потерь осуществляется с помощью стохастического градиентного спуска. 
На каждой итерации параметры обновляются по правилу
\begin{equation}
\theta_{k+1} = \theta_k - \eta \nabla \mathcal{L}_{\mathcal{B}_k}(\theta_k),
\end{equation}
где $\mathcal{B}_k$ — мини-батч размера $B$, а $\nabla \mathcal{L}_{\mathcal{B}_k}$ является стохастической оценкой полного градиента.

В окрестности стационарной точки динамика стохастического градиентного спуска может быть аппроксимирована стохастическим дифференциальным уравнением, в котором уровень шума определяется параметрами обучения. В этом контексте вводится \emph{эффективная температура обучения}
\begin{equation}
T \propto \frac{\eta}{B},
\end{equation}
характеризующая интенсивность стохастических флуктуаций параметров.

Различные значения эффективной температуры соответствуют различным \emph{режимам стохастической оптимизации}, в которых качественно меняется характер динамики SGD, устойчивость решений и исследуемая область параметрического пространства.

Локальная геометрия достигнутых решений описывается кривизной функции потерь в окрестности стационарной точки, характеризуемой спектром матрицы Гессе
\begin{equation}
H(\theta) = \nabla^2 \mathcal{L}(\theta).
\end{equation}
Известно, что многие определения остроты или плоскости минимума не инвариантны к перепараметризациям модели, поэтому в настоящей работе кривизна используется исключительно в виде прокси-оценок для сравнительного анализа в рамках фиксированной архитектуры и параметризации.
